% Part: first-order-logic
% Chapter: proof-systems
% Section: introduction

\documentclass[../../../include/open-logic-section]{subfiles}

\begin{document}

\iftag{FOL}
      {\olfileid{fol}{prf}{int}}
      {\olfileid{pl}{prf}{int}}

\olsection{مقدمه}

منطق‌ها معمولاً هم معناشناسی دارند و هم دستگاهی برای !!a{derivation}.
معناشناسی به مفاهیمی چون صدق، ارضاپذیری، اعتبار و استلزام می‌پردازد.
هدف دستگاه‌های !!{derivation} فراهم‌آوردن روشی صرفاً نحوی برای اثبات
استلزام و اعتبار است. این دستگاه‌ها از آن رو صرفاً نحوی‌اند که
!!{derivation} در چنین دستگاهی شیئی نحوی و متناهی است: معمولاً دنباله‌ای
(یا آرایش متناهیِ دیگری) از !!{sentence}s یا !!{formula}s. دستگاه‌های
خوبِ !!{derivation} این خاصیت را دارند که درستیِ هر دنباله یا آرایشِ
داده‌شده از !!{sentence}s یا !!{formula}s را بتوان به‌طور مکانیکی
بررسی کرد.

ساده‌ترین (و از نظر تاریخی نخستین) دستگاه‌های !!{derivation} برای
منطق مرتبهٔ اول، \emph{اصل‌موضوعی} بودند. دنباله‌ای از !!{formula}s در
چنین دستگاهی !!a{derivation} به‌شمار می‌آید اگر هر !!{formula} آن یا
عضو مجموعه‌ای ثابت از «اصول موضوع» باشد، یا با یکی از شمار ثابتی از
«قواعد استنتاج» از !!{formula}s پیشینِ دنباله به‌دست آید؛ همچنین
می‌توان به‌طور مکانیکی بررسی کرد که آیا !!a{formula} اصل موضوع است و
آیا با یکی از قواعد استنتاج به‌درستی از !!{formula}s دیگر به‌دست
می‌آید. توصیف دستگاه‌های اصل‌موضوعیِ !!{derivation} آسان است و بررسی
فرانظریِ آن‌ها نیز آسان است؛ اما خواندن و فهمیدن !!{derivation}s در
این دستگاه‌ها دشوار است و ساختنشان نیز دشوار.

دستگاه‌های دیگری برای !!{derivation} پدید آمده‌اند تا ساختن
!!{derivation}s یا فهمیدن !!{derivation}s تکمیل‌شده را آسان‌تر کنند.
از نمونه‌های آن‌ها می‌توان استنتاج طبیعی، درخت‌های صدق (که برهان‌های
تابلویی نیز نامیده می‌شوند) و حساب دنباله‌ای را نام برد. برخی
دستگاه‌های !!{derivation} به‌ویژه با هدف مکانیزه‌سازی طراحی شده‌اند؛
برای مثال، پیاده‌سازی روش رزولوشن در نرم‌افزار آسان است (اما فهمیدن
!!{derivation}s آن عملاً ناممکن است). بیشتر این دستگاه‌های دیگر برای
!!{derivation}، !!{derivation}s را به‌جای دنباله، به‌شکل درخت‌هایی از
!!{formula}s
نمایش می‌دهند. ازاین‌رو آسان‌تر می‌توان دید که کدام بخش‌های
!!a{derivation} به کدام بخش‌های دیگر وابسته‌اند.

پس برای منطقی معین، مانند منطق مرتبهٔ اول، دستگاه‌های گوناگونِ
!!{derivation} توضیح‌های متفاوتی به دست می‌دهند از اینکه !!a{sentence}
چه هنگام \emph{قضیه} است و !!a{sentence} چه هنگام از چند جملهٔ دیگر
!!{derivable} است. این کار به هر روشی که انجام شود (با !!{derivation}sی
اصل‌موضوعی، استنتاج‌های طبیعی، !!{derivation}sی دنباله‌ای، درخت‌های صدق
یا ابطال‌های رزولوشنی)، می‌خواهیم این روابط با مفاهیم معناییِ اعتبار و
استلزام منطبق باشند. $\Proves !A$ را برای «$!A$~قضیه است» و
«$\Gamma \Proves !A$» را برای «$!A$ از~$\Gamma$ !!{derivable} است»
می‌نویسیم. $\Proves$ به هر نحوی که تعریف شود، می‌خواهیم با $\Entails$
منطبق باشد؛ یعنی:
\begin{enumerate}
\item $\Proves !A$ اگر و تنها اگر $\Entails !A$
\item $\Gamma \Proves !A$ اگر و تنها اگر $\Gamma \Entails !A$
\end{enumerate}
جهت «فقط اگر» در بالا \emph{صحت} نامیده می‌شود. دستگاه
!!^a{derivation} صحیح است اگر !!{derivability} استلزام (یا اعتبار) را
تضمین کند.
هر دستگاه قابل‌قبولی برای !!{derivation} باید صحیح باشد؛ دستگاه‌های
ناصحیحِ !!{derivation} هیچ فایده‌ای ندارند. به‌هرحال، تمام مقصود از
!!a{derivation} فراهم‌کردن تضمینی نحوی برای اعتبار یا استلزام است.
صحت دستگاه‌های !!{derivation} معرفی‌شده را ثابت خواهیم کرد.

جهت عکس، یعنی «اگر»، نیز مهم است و \emph{تمامیت} نامیده می‌شود.
دستگاه تامِ !!{derivation} آن‌قدر قوی است که هرگاه $!A$ معتبر باشد،
قضیه‌بودنِ $!A$ را نشان دهد، و رابطهٔ $\Gamma \Proves !A$ را هرگاه
$\Gamma \Entails !A$، به دست دهد. اثبات تمامیت دشوارتر است و برخی
منطق‌ها هیچ دستگاه تامی برای !!{derivation} ندارند؛ اما منطق مرتبهٔ
اول چنین دستگاهی دارد. کورت گودل نخستین کسی بود که در رسالهٔ دکتریِ
۱۹۲۹ خود تمامیتِ دستگاهی برای !!a{derivation} در منطق مرتبهٔ اول را
ثابت کرد.

مفهوم دیگری که با دستگاه‌های !!{derivation} پیوند دارد \emph{سازگاری}
است. مجموعه‌ای از !!{sentence}s را ناسازگار می‌نامیم اگر هر چیزی را
بتوان از آن !!{derive} کرد، و در غیر این صورت آن را سازگار می‌نامیم.
ناسازگاری همتای نحویِ ارضاناپذیری است: مجموعه‌های ناسازگارِ
!!{sentence}s نیز، همانند مجموعه‌های ارضاناپذیر، نظریه‌های خوبی نیستند
و نقصی بنیادی دارند. مجموعه‌های سازگارِ !!{sentence}s ممکن است صادق
یا سودمند نباشند، اما دست‌کم از این آستانهٔ کمینهٔ سودمندی منطقی
می‌گذرند. تعریف دقیق سازگاری ممکن است در دستگاه‌های گوناگونِ
!!{derivation} برای مجموعه‌های !!{sentence}s فرق کند؛ اما
همانند~$\Proves$،
می‌خواهیم سازگاری با همتای معناییِ خود، یعنی ارضاپذیری، منطبق باشد.
می‌خواهیم همواره $\Gamma$ سازگار باشد اگر و تنها اگر ارضاپذیر باشد.
در اینجا جهت «فقط اگر» همان تمامیت است (سازگاری، ارضاپذیری را تضمین
می‌کند) و جهت «اگر» همان صحت است (ارضاپذیری، سازگاری را تضمین می‌کند).
در واقع، برای منطق کلاسیک مرتبهٔ اول، دو صورتِ صحت و تمامیت هم‌ارزند.

\end{document}
